
\begin{lemma}[Two-site support after refinement]
    \label{lem:mpdo_vertical_sector_t_image_preservation}
    \lean{MPOTensor.precompressionVerticalSectorT_image_preservation}
    \leanok
    \uses{def:mpdo_vertical_sector_precompression,
      lem:mpdo_normalized_retained_vertical_sector_embedding,
      lem:mpdo_retained_vertical_weighted_contraction,
      lem:mpdo_vertical_bond_contraction_reconstruction,
      thm:mpdo_one_site_closure_two_site_blocking}
    For a bond matrix $X$, set
    \[
        V_1(X)=\bigoplus_\alpha m_\alpha M_\alpha(X).
    \]
    Suppose that both vertical canonical forms satisfy their exact
    reconstruction identities, the one-site multiplicity spaces are nonzero,
    their diagonal weights are positive,
    $T(\mathcal M^{(1)}(X))=\mathcal M^{(2)}(X)$, and
    $U_2U_2^\dagger=1$, then
    \[
        U_2^\dagger U_2 Z_T(V_1(X))=Z_T(V_1(X)).
    \]
    This is the refinement support identity used in
    \cite[Appendix~C.4, lines~1955--1980]{Cirac2016MPDO_arXiv}.
\end{lemma}

\begin{proof}\leanok
    Let $\mathcal R^{-1}$ relabel a matrix indexed by two one-site indices as
    a matrix indexed by the blocked two-site index.  Then
    \[
        Z_T(V_1(X))
        =\mathcal R^{-1}\!\left[
          T\!\left(U_1^\dagger B_1(X)U_1\right)\right]
        =\mathcal R^{-1}\!\left[
          T\!\left(\mathcal M^{(1)}(X)\right)\right]
        =\mathcal R^{-1}\!\left[\mathcal M^{(2)}(X)\right]
        =U_2^\dagger B_2(X)U_2.
    \]
    Here the first equality is the normalized embedding, the second and last
    are the two reconstruction identities, and the middle equality is the
    refinement identity.  The equality before the last is precisely the
    blocked-index relabelling theorem.
    Therefore
    \[
        U_2^\dagger U_2Z_T(V_1(X))
        =U_2^\dagger(U_2U_2^\dagger)B_2(X)U_2
        =Z_T(V_1(X)).
    \]
\end{proof}

\begin{lemma}[One-site support after coarse-graining]
    \label{lem:mpdo_vertical_sector_s_image_preservation}
    \lean{MPOTensor.precompressionVerticalSectorS_image_preservation}
    \leanok
    \uses{def:mpdo_vertical_sector_precompression,
      lem:mpdo_normalized_retained_vertical_sector_embedding,
      lem:mpdo_retained_vertical_weighted_contraction,
      lem:mpdo_vertical_bond_contraction_reconstruction,
      thm:mpdo_one_site_closure_two_site_blocking}
    For a bond matrix $X$, set
    \[
        V_2(X)=\bigoplus_\beta n_\beta A_\beta(X).
    \]
    Suppose that both vertical canonical forms satisfy their exact
    reconstruction identities, the two-site multiplicity spaces are nonzero,
    their diagonal weights are positive,
    $S(\mathcal M^{(2)}(X))=\mathcal M^{(1)}(X)$, and
    $U_1U_1^\dagger=1$, then
    \[
        U_1^\dagger U_1 Z_S(V_2(X))=Z_S(V_2(X)).
    \]
    This is the coarse-graining support identity used in
    \cite[Appendix~C.4, lines~1955--1980]{Cirac2016MPDO_arXiv}.
\end{lemma}

\begin{proof}\leanok
    Let $\mathcal R$ decode a matrix indexed by the blocked two-site index as a
    matrix indexed by two one-site indices.  Then
    \[
        Z_S(V_2(X))
        =S\!\left(\mathcal R\!\left[
          U_2^\dagger B_2(X)U_2\right]\right)
        =S\!\left(\mathcal M^{(2)}(X)\right)
        =\mathcal M^{(1)}(X)
        =U_1^\dagger B_1(X)U_1.
    \]
    Here the first equality is the normalized embedding, the next equality is
    the blocked-index relabelling theorem, the middle equality is the
    coarse-graining identity, and the last is the one-site reconstruction.
    Therefore
    \[
        U_1^\dagger U_1Z_S(V_2(X))
        =U_1^\dagger(U_1U_1^\dagger)B_1(X)U_1
        =Z_S(V_2(X)).
    \]
\end{proof}

\begin{lemma}[Trace preservation on one-site canonical-form contractions]
    \label{lem:mpdo_vertical_sector_t_source_generated_trace}
    \lean{MPOTensor.transportedVerticalSectorT_trace_of_source_generated}
    \leanok
    \uses{lem:mpdo_vertical_sector_t_image_preservation,
      lem:mpdo_vertical_sector_precompression_positive_trace}
    Under the hypotheses of Lemma~\ref{lem:mpdo_vertical_sector_t_image_preservation},
    suppose in addition that $U_1U_1^\dagger=1$ and that $T$ is completely
    positive and trace preserving.  Then
    \[
        \Trsec(\widetilde T(V_1(X)))=\Trsec(V_1(X)).
    \]
\end{lemma}

\begin{proof}\leanok
    Cyclicity of trace and the preceding support identity give
    \[
        \Trsec(\widetilde T(V_1(X)))
        =\tr(U_2^\dagger U_2Z_T(V_1(X)))
        =\tr(Z_T(V_1(X)))
        =\Trsec(V_1(X)).
    \]
\end{proof}

\begin{lemma}[Trace preservation on two-site canonical-form contractions]
    \label{lem:mpdo_vertical_sector_s_source_generated_trace}
    \lean{MPOTensor.transportedVerticalSectorS_trace_of_source_generated}
    \leanok
    \uses{lem:mpdo_vertical_sector_s_image_preservation,
      lem:mpdo_vertical_sector_precompression_positive_trace}
    Under the hypotheses of Lemma~\ref{lem:mpdo_vertical_sector_s_image_preservation},
    suppose in addition that $U_2U_2^\dagger=1$ and that $S$ is completely
    positive and trace preserving.  Then
    \[
        \Trsec(\widetilde S(V_2(X)))=\Trsec(V_2(X)).
    \]
\end{lemma}

\begin{proof}\leanok
    Cyclicity of trace and the preceding support identity give
    \[
        \Trsec(\widetilde S(V_2(X)))
        =\tr(U_1^\dagger U_1Z_S(V_2(X)))
        =\tr(Z_S(V_2(X)))
        =\Trsec(V_2(X)).
    \]
\end{proof}

\begin{theorem}[Action of the normalized vertical-sector maps]
    \label{thm:mpdo_transported_vertical_sector_maps}
    \lean{MPOTensor.transportedVerticalSectorT_trace_smul}
    \lean{MPOTensor.transportedVerticalSectorS_trace_smul}
    \lean{MPOTensor.transportedVerticalSectorT_contractBondMatrix_trace_smul}
    \lean{MPOTensor.transportedVerticalSectorS_contractBondMatrix_trace_smul}
    \leanok
    \uses{def:mpdo_transported_vertical_sector_maps,
      lem:mpdo_normalized_retained_vertical_sector_embedding,
      lem:mpdo_retained_vertical_sector_partial_trace,
      thm:mpdo_vertical_coordinate_map_transport}
    Suppose that every multiplicity space for $\mu_\alpha$ and $\nu_\beta$
    is nonzero, every diagonal entry of these two matrices is positive, and
    \[
        T_{\mathrm v}\!\left(
          \bigoplus_\alpha\mu_\alpha\otimes X_\alpha\right)
        =\bigoplus_\beta\nu_\beta\otimes Y_\beta,
        \qquad
        S_{\mathrm v}\!\left(
          \bigoplus_\beta\nu_\beta\otimes Y_\beta\right)
        =\bigoplus_\alpha\mu_\alpha\otimes X_\alpha.
    \]
    Writing $m_\alpha=\tr(\mu_\alpha)$ and
    $n_\beta=\tr(\nu_\beta)$, one has
    \[
        \widetilde T\!\left(
          \bigoplus_\alpha m_\alpha X_\alpha\right)
        =\bigoplus_\beta n_\beta Y_\beta,
        \qquad
        \widetilde S\!\left(
          \bigoplus_\beta n_\beta Y_\beta\right)
        =\bigoplus_\alpha m_\alpha X_\alpha.
    \]
    In particular, take $X_\alpha=M_\alpha(X)$ and
    $Y_\beta=A_\beta(X)$ from the contractions of the one-site and two-site
    vertical canonical forms.  The physical closure identities for $T$ and
    $S$ then imply the two hypotheses above, so the same conclusions hold for
    these source-generated sector operators.
\end{theorem}

\begin{proof}\leanok
    Since $\widetilde T=\widetilde R_2T_{\mathrm v}R_1$, one has
    \[
    \begin{aligned}
        \widetilde T\!\left(\bigoplus_\alpha m_\alpha X_\alpha\right)
        &=\widetilde R_2\!\left(T_{\mathrm v}\!\left(
          \bigoplus_\alpha\mu_\alpha\otimes X_\alpha\right)\right) \\
        &=\widetilde R_2\!\left(
          \bigoplus_\beta\nu_\beta\otimes Y_\beta\right)
        =\bigoplus_\beta n_\beta Y_\beta.
    \end{aligned}
    \]
    Likewise, $\widetilde S=\widetilde R_1S_{\mathrm v}R_2$ gives
    \[
    \begin{aligned}
        \widetilde S\!\left(\bigoplus_\beta n_\beta Y_\beta\right)
        &=\widetilde R_1\!\left(S_{\mathrm v}\!\left(
          \bigoplus_\beta\nu_\beta\otimes Y_\beta\right)\right) \\
        &=\widetilde R_1\!\left(
          \bigoplus_\alpha\mu_\alpha\otimes X_\alpha\right)
        =\bigoplus_\alpha m_\alpha X_\alpha.
    \end{aligned}
    \]
\end{proof}

\begin{lemma}[Fixed contraction families of the transported composites]
    \label{lem:mpdo_transported_vertical_sector_fixed_generators}
    \lean{MPOTensor.transportedVerticalSectorS_comp_T_fixed_contractBondMatrix_trace_smul}
    \lean{MPOTensor.transportedVerticalSectorT_comp_S_fixed_contractBondMatrix_trace_smul}
    \leanok
    \uses{thm:mpdo_transported_vertical_sector_maps}
    For every bond matrix $X$, set
    \[
        V_1(X)=\bigoplus_\alpha m_\alpha M_\alpha(X),
        \qquad
        V_2(X)=\bigoplus_\beta n_\beta A_\beta(X).
    \]
    Under the hypotheses of the preceding theorem, the transported
    composites satisfy
    \[
        (\widetilde S\circ\widetilde T)(V_1(X))=V_1(X),
        \qquad
        (\widetilde T\circ\widetilde S)(V_2(X))=V_2(X).
    \]
    Thus the one-site and two-site contraction families lie in the respective
    fixed-point spaces, as in Appendix~C.4, lines 1974--1980.  These equations
    do not assert that either composite is the identity on the whole sector
    algebra; that conclusion requires the fixed-point structure argument of
    Appendix~C.4, lines 1980--1993.
\end{lemma}

\begin{proof}\leanok
    The preceding theorem gives
    $\widetilde T(V_1(X))=V_2(X)$ and
    $\widetilde S(V_2(X))=V_1(X)$.  Hence
    \[
        (\widetilde S\circ\widetilde T)(V_1(X))
        =\widetilde S(V_2(X))=V_1(X),
        \qquad
        (\widetilde T\circ\widetilde S)(V_2(X))
        =\widetilde T(V_1(X))=V_2(X).
    \]
\end{proof}

\begin{theorem}[Faithful fixed point from fixed product generators]
    \label{thm:positive_map_faithful_fixed_point_from_product_span}
    \lean{IsPositiveMap.exists_posDef_fixedPoint_of_traceNonincreasing_of_fixed_product_span}
    \leanok
    Let $\mathcal M$ be a finite full matrix algebra over $\mathbb C$, and let
    $F:\mathcal M\to\mathcal M$ be positive and trace nonincreasing.  Suppose
    that a family $(V_j)_{j\in J}$ consists of fixed points of $F$ and that,
    for some $L>0$, the identity belongs to the linear span of the products
    \[
      V_{j_1}\cdots V_{j_L}
    \]
    in $\mathcal M$.  Then $F$ has a positive-definite fixed point.
    This is the local full-matrix intermediate implication used in
    \cite[Appendix~C.4, lines~1980--1995]{Cirac2016MPDO_arXiv}; it is not
    stated there as a separate theorem.
\end{theorem}

\begin{proof}\leanok
    \uses{thm:positive_trace_nonincreasing_bounded_orbits,
      thm:maximalSupport_fixedPoint_of_bounded_orbits}
    Boundedness of the forward orbits gives the mean-ergodic projection $P$.
    Set $\rho=P(\mathbf 1)$ and let $Q$ be the support projection of $\rho$.
    For every $j$, the maximal-support property of $\rho$ gives
    \[
      QV_jQ=V_j.
    \]
    Since $Q$ is a projection, this implies $QV_j=V_j=V_jQ$.  By induction,
    \[
      Q\cdot V_{j_1}\cdots V_{j_k}\cdot Q
        =V_{j_1}\cdots V_{j_k}
      \qquad (k\geq1).
    \]
    The product-span hypothesis and linearity then give
    $Q\mathbf 1Q=\mathbf 1$.  Since $Q^2=Q$, one has
    \[
      Q=Q^2=Q\mathbf 1Q=\mathbf 1.
    \]
    Thus $\rho$ is positive definite, and the construction of $P$ gives
    $F(\rho)=\rho$.
\end{proof}

\begin{theorem}[Trace preservation from fixed product generators]
    \label{thm:positive_map_trace_preserving_from_fixed_product_span}
    \lean{IsPositiveMap.tracePreserving_of_traceNonincreasing_of_fixed_product_span}
    \leanok
    Let $\mathcal M$ be a finite full matrix algebra over $\mathbb C$, and let
    $F:\mathcal M\to\mathcal M$ be positive and trace nonincreasing.  Suppose
    that $F(V_j)=V_j$ for every $j\in J$.  Put
    $\mathcal V=\spn_{\C}\{V_j:j\in J\}$ and, in the notation of
    \cite[Appendix~C.4]{Cirac2016MPDO_arXiv}, define
    \[
      C_L(\mathcal V)
        =\spn_{\C}\{v_1\cdots v_L:v_1,\ldots,v_L\in\mathcal V\}.
    \]
    If $\mathbf 1\in C_L(\mathcal V)$ for some $L>0$, then $F$ is trace
    preserving.
    % The product-span hypothesis is the fixed-generator structure appearing in
    % \cite[Appendix~C.4, lines~1974--1995]{Cirac2016MPDO_arXiv}.  Trace
    % preservation is derived here from trace nonincrease and the resulting
    % faithful fixed point, not from the cited passage.
\end{theorem}

\begin{proof}\leanok
    \uses{thm:positive_map_faithful_fixed_point_from_product_span,
      thm:trace_pairing_adjoint_positive,
      thm:psd_trace_pairing_self_dual,
      lem:posDef_weighted_trace_faithful,
      thm:trace_preserving_iff_trace_adjoint_unital}
    Multilinearity gives
    \[
      C_L(\mathcal V)
        =\spn_{\C}\{V_{j_1}\cdots V_{j_L}:j_1,\ldots,j_L\in J\}.
    \]
    By Theorem~\ref{thm:positive_map_faithful_fixed_point_from_product_span},
    there is a positive-definite matrix $\rho$ such that $F(\rho)=\rho$.
    Let $F^*$ be the trace-pairing adjoint and set
    \[
      \Delta=\mathbf 1-F^*(\mathbf 1).
    \]
    For every positive semidefinite $X$, trace nonincrease gives
    \[
      \tr(\Delta X)
        =\tr(X)-\tr(F(X))\geq0.
    \]
    Positivity of $F^*$ and self-duality of the positive semidefinite cone
    therefore give $\Delta\geq0$.  The fixed-point equation gives
    $\tr(\rho\Delta)=0$, so positive definiteness of $\rho$
    forces $\Delta=0$.  Thus $F^*(\mathbf 1)=\mathbf 1$, which is equivalent
    to trace preservation.
\end{proof}

\begin{theorem}[Trace transfer through the canonical full-matrix extension]
    \label{thm:direct_sum_extension_trace_transfer}
    \lean{Matrix.IsTraceNonincreasingBetweenDirectSums.directSumExtension}
    \lean{Matrix.IsTracePreservingMap.isTracePreservingDirectSumMap_of_extension}
    \leanok
    \uses{def:direct_sum_full_matrix_extension}
    Let
    $\mathcal A=\bigoplus_{\alpha\in I}\mathcal M_{d_\alpha\times d_\alpha}$,
    let $\iota:\mathcal A\to\End(\bigoplus_\alpha
    \mathbb C^{d_\alpha})$ be the block-diagonal embedding, and let $\pi$ be
    diagonal compression.  For a linear map $F:\mathcal A\to\mathcal A$, set
    $\widehat F=\iota\circ F\circ\pi$.  Then:
    \begin{enumerate}
      \item if $F$ is trace nonincreasing on positive elements of $\mathcal A$,
        then $\widehat F$ is trace nonincreasing;
      \item if $\widehat F$ is trace preserving, then $F$ preserves the total
        trace on $\mathcal A$.
    \end{enumerate}
\end{theorem}

\begin{proof}\leanok
    If $Y\geq0$, then $\pi(Y)\geq0$, and the trace identities for diagonal
    compression and block-diagonal embedding give
    \[
      \tr(\widehat F(Y))
        =\tr_{\mathcal A}(F(\pi(Y)))
        \leq\tr_{\mathcal A}(\pi(Y))
        =\tr(Y).
    \]
    This proves the first assertion.  For $X\in\mathcal A$, the identity
    $\pi\circ\iota=\id_{\mathcal A}$ gives
    \[
      \tr_{\mathcal A}(F(X))
        =\tr(\widehat F(\iota(X)))
        =\tr(\iota(X))
        =\tr_{\mathcal A}(X),
    \]
    whenever $\widehat F$ is trace preserving.
\end{proof}

\begin{theorem}[Trace preservation on a finite sum from fixed products]
    \label{thm:positive_direct_sum_map_trace_preserving_from_fixed_product_span}
    \lean{Matrix.IsPositiveDirectSumMap.tracePreserving_of_traceNonincreasing_of_fixed_product_span}
    \leanok
    Let
    $\mathcal A=\bigoplus_{\alpha\in I}\mathcal M_{d_\alpha\times d_\alpha}$
    be a finite sum of full matrix algebras, with total trace
    \[
      \tr_{\mathcal A}(X)
        =\sum_{\alpha\in I}\tr(X_\alpha).
    \]
    Let $F:\mathcal A\to\mathcal A$ be positive and trace nonincreasing.
    Suppose that $F(V_j)=V_j$ for every $j\in J$.  With
    $\mathcal V=\spn_{\C}\{V_j:j\in J\}$, put
    \[
      C_L(\mathcal V)
        =\spn_{\C}\{v_1\cdots v_L:v_1,\ldots,v_L\in\mathcal V\}.
    \]
    If $\mathbf 1_{\mathcal A}\in C_L(\mathcal V)$ for some $L>0$, then $F$
    preserves the total trace on $\mathcal A$.
    % This is the finite-sector form of the preceding full-matrix trace argument.
    % The source passage
    % \cite[Appendix~C.4, lines~1974--1995]{Cirac2016MPDO_arXiv} supplies the
    % product-span structure to which that argument is applied.
\end{theorem}

\begin{proof}\leanok
    \uses{thm:positive_map_trace_preserving_from_fixed_product_span,
      thm:direct_sum_full_matrix_extension_compatibility,
      thm:direct_sum_extension_trace_transfer}
    Multilinearity first gives
    \[
      C_L(\mathcal V)
        =\spn_{\C}\{V_{j_1}\cdots V_{j_L}:j_1,\ldots,j_L\in J\}.
    \]
    Let $\iota:\mathcal A\to\End(\bigoplus_\alpha
    \mathbb C^{d_\alpha})$ be the block-diagonal embedding and let
    $\widehat F$ be the canonical full-matrix extension.  These maps satisfy
    \[
      \iota(XY)=\iota(X)\iota(Y),
      \qquad \iota(\mathbf 1_{\mathcal A})=\mathbf 1,
      \qquad \widehat F(\iota(X))=\iota(F(X)).
    \]
    The extension is positive and trace nonincreasing.  The matrices
    $\iota(V_j)$ are fixed by $\widehat F$, and the displayed identities
    carry the product-span hypothesis to the full matrix algebra.
    Theorem~\ref{thm:positive_map_trace_preserving_from_fixed_product_span}
    makes $\widehat F$ trace preserving.  Restricting to block-diagonal
    matrices proves preservation of $\tr_{\mathcal A}$ by $F$.
\end{proof}

\begin{theorem}[Trace sandwich on finite sums of matrix algebras]
    \label{thm:direct_sum_trace_preservation_sandwich}
    \lean{Matrix.isTracePreservingBetweenDirectSums_of_comp_of_traceNonincreasing}
    \leanok
    Let $\mathcal A$ and $\mathcal B$ be finite sums of full matrix algebras,
    and let $T:\mathcal A\to\mathcal B$ and $S:\mathcal B\to\mathcal A$ be
    linear maps.  Suppose that $T$ is positive, both maps are trace
    nonincreasing, and $S\circ T$ is trace preserving.  Then $T$ is trace
    preserving.
\end{theorem}

\begin{proof}\leanok
    For every positive $X\in\mathcal A$, positivity of $T$ and the two
    trace-nonincreasing inequalities give
    \[
      \tr_{\mathcal A}(X)
        =\tr_{\mathcal A}(S(T(X)))
        \leq\tr_{\mathcal B}(T(X))
        \leq\tr_{\mathcal A}(X).
    \]
    Hence the middle term equals the endpoints.  Equality extends from the
    positive cone to Hermitian matrices by positive and negative parts.  For
    arbitrary $X$, put
    \[
      H_1=X+X^\dagger,
      \qquad H_2=\mathrm{i}(X-X^\dagger),
      \qquad X=\frac12(H_1-\mathrm{i}H_2).
    \]
    Both $H_1$ and $H_2$ are Hermitian, so linearity proves trace preservation
    for every $X$.
\end{proof}

\begin{theorem}[Trace preservation of the transported vertical-sector maps]
    \label{thm:mpdo_transported_vertical_sector_composites_trace_preserving}
    \lean{MPOTensor.transportedVerticalSector_composites_tracePreserving}
    \leanok
    Suppose that every one-site and two-site multiplicity space is nonzero
    and every corresponding diagonal weight is positive.  Let
    $\widetilde M^{(1)}$ and $\widetilde M^{(2)}$ be the vertical tensors of
    $M$ and its two-site blocking, and let $B_1$ and $B_2$ be their assembled
    one-site and two-site canonical forms.  Assume that both canonical forms
    are bases of normal tensors and that, for every vertical letter $\ell$,
    \[
      U_i\widetilde M^{(i)}_\ell U_i^\dagger=(B_i)_\ell,
      \qquad
      \widetilde M^{(i)}_\ell
        =U_i^\dagger(B_i)_\ell U_i,
      \qquad
      U_iU_i^\dagger=\mathbf 1
      \quad (i=1,2).
    \]
    Assume in addition that the physical maps $T$ and $S$ are completely
    positive and trace preserving and that, for every horizontal bond matrix
    $Z$, the two physical-coordinate identities hold:
    \[
      T(\mathcal M^{(1)}(Z))=\mathcal M^{(2)}(Z),
      \qquad
      S(\mathcal M^{(2)}(Z))=\mathcal M^{(1)}(Z).
    \]
    Then each transported map and both square composites preserve the
    appropriate total sector trace.
    For the individual maps,
    \[
      \sum_\beta\tr(\widetilde T(X)_\beta)
        =\sum_\alpha\tr(X_\alpha),
      \qquad
      \sum_\alpha\tr(\widetilde S(Y)_\alpha)
        =\sum_\beta\tr(Y_\beta).
    \]
    For the square composites,
    \[
      \sum_\alpha\tr
        \bigl((\widetilde S\widetilde T)(X)_\alpha\bigr)
        =\sum_\alpha\tr(X_\alpha),
      \qquad
      \sum_\beta\tr
        \bigl((\widetilde T\widetilde S)(Y)_\beta\bigr)
        =\sum_\beta\tr(Y_\beta).
    \]
    Neither of the stronger range inclusions
    \[
      T(R_1(\mathcal A_1))\subseteq R_2(\mathcal A_2),
      \qquad
      S(R_2(\mathcal A_2))\subseteq R_1(\mathcal A_1)
    \]
    is assumed.
\end{theorem}

\begin{proof}\leanok
    \uses{thm:positive_direct_sum_map_trace_preserving_from_fixed_product_span,
      thm:direct_sum_trace_preservation_sandwich,
      lem:mpdo_transported_vertical_sector_fixed_generators,
      thm:mpdo_vertical_bond_contractions_generate_sector_algebra,
      thm:mpdo_transported_vertical_sector_trace_loss}
    Each transported map is positive and trace nonincreasing, and hence so are
    the square composites $\widetilde S\widetilde T$ and
    $\widetilde T\widetilde S$.  The fixed-generator lemma and the
    product-spanning theorem give the hypotheses of
    Theorem~\ref{thm:positive_direct_sum_map_trace_preserving_from_fixed_product_span}
    for both composites.  Thus both preserve the corresponding total trace.

    For positive $X\in\mathcal A_1$, one has
    \[
      \tr_{\mathcal A_1}(X)
        =\tr_{\mathcal A_1}
          ((\widetilde S\widetilde T)(X))
        \leq\tr_{\mathcal A_2}(\widetilde T(X))
        \leq\tr_{\mathcal A_1}(X).
    \]
    Theorem~\ref{thm:direct_sum_trace_preservation_sandwich} therefore makes
    $\widetilde T$ trace preserving.  Applying the same theorem to
    $\widetilde S$ and $\widetilde T\widetilde S$ proves the assertion for
    $\widetilde S$.
\end{proof}

\begin{lemma}[Complete positivity of refinement in vertical coordinates]
    \label{lem:mpdo_transport_t_vertical_coordinates_cp}
    \lean{MPOTensor.transportTToVerticalCoordinates_isKrausCP}
    \leanok
    \uses{def:mpdo_vertical_coordinate_map_transport}
    If the physical refinement map $T$ is completely positive, then the
    transported map $T_{\mathrm v}$ in vertical canonical coordinates is
    completely positive.
\end{lemma}

\begin{proof}\leanok
    By Definition~\ref{def:mpdo_vertical_coordinate_map_transport},
    \[
      T_{\mathrm v}(Z)
        =U_2J\!\left(T(U_1^\dagger ZU_1)\right)U_2^\dagger.
    \]
    The two single-operator conjugations and the reindexing $J$ are completely
    positive, so the assertion follows by composition.
\end{proof}

\begin{lemma}[Complete positivity of coarse-graining in vertical coordinates]
    \label{lem:mpdo_transport_s_vertical_coordinates_cp}
    \lean{MPOTensor.transportSToVerticalCoordinates_isKrausCP}
    \leanok
    \uses{def:mpdo_vertical_coordinate_map_transport}
    If the physical coarse-graining map $S$ is completely positive, then the
    transported map $S_{\mathrm v}$ in vertical canonical coordinates is
    completely positive.
\end{lemma}

\begin{proof}\leanok
    By Definition~\ref{def:mpdo_vertical_coordinate_map_transport},
    \[
      S_{\mathrm v}(W)
        =U_1S\!\left(J^{-1}(U_2^\dagger WU_2)\right)U_1^\dagger.
    \]
    Again, single-operator conjugation, reindexing, and composition preserve
    complete positivity.
\end{proof}

\begin{lemma}[Retained-coordinate block-diagonal inclusion]
    \label{lem:mpdo_vertical_sector_block_diagonal_reindex}
    \lean{MPOTensor.verticalSectorBlockDiagonal_eq_equivReindexMap_comp_embedding}
    \leanok
    \uses{def:mpdo_retained_vertical_sector_coordinates}
    Let $\iota$ be block-diagonal inclusion of the weighted vertical-sector
    block family into its full matrix algebra, and let $W$ be the canonical
    reindexing to retained physical coordinates.  Then the retained-coordinate
    inclusion is
    \[
      \iota_{\mathrm v}=W\circ\iota.
    \]
\end{lemma}

\begin{proof}\leanok
    This is the defining relation between the direct-sum and retained
    coordinates.
\end{proof}

\begin{lemma}[Retained-coordinate diagonal-block extraction]
    \label{lem:mpdo_vertical_sector_blocks_reindex}
    \lean{MPOTensor.verticalSectorBlocks_eq_compression_comp_equivReindexMap}
    \leanok
    \uses{def:mpdo_retained_vertical_sector_coordinates}
    Let $\pi$ be diagonal-block extraction to the weighted vertical-sector
    block family, and let $W$ be the canonical reindexing to retained physical
    coordinates.  Then the retained-coordinate extraction is
    \[
      \pi_{\mathrm v}=\pi\circ W^{-1}.
    \]
\end{lemma}

\begin{proof}\leanok
    This is the inverse defining relation between the retained and direct-sum
    coordinates.
\end{proof}

\begin{lemma}[Full-matrix factorization of transported refinement]
    \label{lem:mpdo_transported_vertical_sector_t_extension}
    \lean{MPOTensor.directSumMapExtension_transportedVerticalSectorT}
    \leanok
    \uses{def:mpdo_transported_vertical_sector_maps,
      lem:mpdo_vertical_sector_block_diagonal_reindex,
      lem:mpdo_vertical_sector_blocks_reindex}
    Write $\widehat R_1$ for normalized state preparation,
    $\widehat{\widetilde R}_2$ for controlled partial trace, and $W_i$ for
    retained-coordinate reindexing.  The full-matrix extension of the
    transported refinement map is
    \[
      \widehat{\widetilde T}
        =\widehat{\widetilde R}_2\circ W_2^{-1}\circ
          T_{\mathrm v}\circ W_1\circ\widehat R_1.
    \]
\end{lemma}

\begin{proof}\leanok
    Substitute the preceding inclusion and extraction identities into
    $\widetilde T=\widetilde R_2T_{\mathrm v}R_1$.
\end{proof}

\begin{lemma}[Full-matrix factorization of transported coarse-graining]
    \label{lem:mpdo_transported_vertical_sector_s_extension}
    \lean{MPOTensor.directSumMapExtension_transportedVerticalSectorS}
    \leanok
    \uses{def:mpdo_transported_vertical_sector_maps,
      lem:mpdo_vertical_sector_block_diagonal_reindex,
      lem:mpdo_vertical_sector_blocks_reindex}
    Write $\widehat R_2$ for normalized state preparation,
    $\widehat{\widetilde R}_1$ for controlled partial trace, and $W_i$ for
    retained-coordinate reindexing.  The full-matrix extension of the
    transported coarse-graining map is
    \[
      \widehat{\widetilde S}
        =\widehat{\widetilde R}_1\circ W_1^{-1}\circ
          S_{\mathrm v}\circ W_2\circ\widehat R_2.
    \]
\end{lemma}

\begin{proof}\leanok
    Substitute the preceding inclusion and extraction identities into
    $\widetilde S=\widetilde R_1S_{\mathrm v}R_2$.
\end{proof}

\begin{theorem}[Complete positivity and adjoint Schwarz for the transported maps]
    \label{thm:mpdo_transported_vertical_sector_composites_schwarz}
    \lean{MPOTensor.transportedVerticalSectorT_isKrausDirectSumMap,
      MPOTensor.transportedVerticalSectorS_isKrausDirectSumMap,
      MPOTensor.transportedVerticalSector_composites_traceAdjointSchwarz}
    \leanok
    \uses{def:mpdo_transported_vertical_sector_maps,
      lem:mpdo_transport_t_vertical_coordinates_cp,
      lem:mpdo_transport_s_vertical_coordinates_cp,
      lem:mpdo_transported_vertical_sector_t_extension,
      lem:mpdo_transported_vertical_sector_s_extension,
      thm:mpdo_vertical_sector_partial_trace_kraus,
      lem:direct_sum_kraus_composition_schwarz,
      thm:mpdo_transported_vertical_sector_composites_trace_preserving}
    Under the hypotheses of
    Theorem~\ref{thm:mpdo_transported_vertical_sector_composites_trace_preserving},
    the maps
    \[
      \widetilde T:\mathcal A_1\longrightarrow\mathcal A_2,
      \qquad
      \widetilde S:\mathcal A_2\longrightarrow\mathcal A_1
    \]
    are completely positive.  The trace adjoints of the two square composites
    satisfy the Schwarz inequality in every simple summand.
\end{theorem}

\begin{proof}\leanok
    Let $\iota_i$ and $\pi_i$ be block-diagonal inclusion and diagonal-block
    extraction, respectively, and let $W_i$ denote the canonical reindexing
    from the weighted vertical-sector coordinates to the retained physical
    coordinates.  The full-matrix extension of the transported refinement
    map has the factorization
    \[
      \widehat{\widetilde T}
        =\widehat{\widetilde R}_2\circ W_2^{-1}\circ
          T_{\mathrm v}\circ W_1\circ\widehat R_1.
    \]
    Here $\widehat R_1$ is normalized state preparation,
    $\widehat{\widetilde R}_2$ is controlled partial trace, and
    \[
      T_{\mathrm v}(Z)
        =U_2J\!\left(T(U_1^\dagger ZU_1)\right)U_2^\dagger.
    \]
    Every factor is completely positive.  The analogous identities
    \[
      S_{\mathrm v}(W)
        =U_1S\!\left(J^{-1}(U_2^\dagger WU_2)\right)U_1^\dagger,
      \qquad
      \widehat{\widetilde S}
        =\widehat{\widetilde R}_1\circ W_1^{-1}\circ
          S_{\mathrm v}\circ W_2\circ\widehat R_2
    \]
    prove complete positivity of $\widetilde S$.  These factorizations are
    exact on the full matrix algebras: $\pi_iW_i^{-1}$ is precisely extraction
    of the retained diagonal sector blocks.  Thus no range inclusion is used.

    Composition gives complete positivity of
    $\widetilde S\widetilde T$ and $\widetilde T\widetilde S$.
    The preceding theorem gives trace preservation of these two composites.
    Their trace adjoints are therefore unital completely positive maps, and
    the Kadison--Schwarz inequality applies in each simple summand.
\end{proof}

\begin{theorem}[Transported vertical-sector composites are identities]
    \label{thm:mpdo_transported_vertical_sector_composites_identity}
    \lean{MPOTensor.transportedVerticalSector_composites_eq_id}
    \leanok
    \uses{thm:mpdo_transported_vertical_sector_composites_trace_preserving,
      thm:mpdo_transported_vertical_sector_composites_schwarz,
      lem:mpdo_transported_vertical_sector_fixed_generators,
      thm:mpdo_vertical_bond_contractions_generate_sector_algebra,
      thm:mpdo_fixed_contractions_trace_adjoint_schwarz_identity}
    Under precisely the hypotheses of
    Theorem~\ref{thm:mpdo_transported_vertical_sector_composites_trace_preserving},
    the transported maps satisfy
    \begin{align}
      \widetilde S\circ\widetilde T
        &=\id_{\mathcal A_1},
        \label{eq:mpdo_vsector_ST_id}\\
      \widetilde T\circ\widetilde S
        &=\id_{\mathcal A_2}.
        \label{eq:mpdo_vsector_TS_id}
    \end{align}
    Thus $\widetilde T$ and $\widetilde S$ are mutually inverse.  This is the
    identity-composition conclusion of
    \cite[Appendix~C.4, lines~1974--1995]{Cirac2016MPDO_arXiv}.
\end{theorem}

\begin{proof}\leanok
    For each square composite $F$, positivity follows from positivity of the
    two transported maps.  Theorems~
    \ref{thm:mpdo_transported_vertical_sector_composites_trace_preserving}
    and~\ref{thm:mpdo_transported_vertical_sector_composites_schwarz} give
    trace preservation and the Schwarz inequality for $F^*$, respectively.
    Lemma~\ref{lem:mpdo_transported_vertical_sector_fixed_generators} gives
    \[
      F(V_i(X))=V_i(X)
    \]
    for every bond matrix $X$.  Theorem~
    \ref{thm:mpdo_vertical_bond_contractions_generate_sector_algebra} gives a
    positive integer $L_i$ such that
    \[
      C_{L_i}(V_i)=\mathcal A_i.
    \]
    Applying Theorem~
    \ref{thm:mpdo_fixed_contractions_trace_adjoint_schwarz_identity} first to
    $F=\widetilde S\circ\widetilde T$ and then to
    $F=\widetilde T\circ\widetilde S$ proves
    equations~(\ref{eq:mpdo_vsector_ST_id}) and
    (\ref{eq:mpdo_vsector_TS_id}).
\end{proof}
